% !TEX root = ../main.tex
\section{Model basics}\label{sec:model}
\subsection{RPV-MSSM with a light bino}
In the R-parity-violating MSSM, the usual MSSM superpotential is extended by the following terms
\cite{Hall:1983id, Allanach:2003eb}:
\begin{equation}
    W_\text{RPV} = \kappa_i L_i H_u + \frac{1}{2} \lambda_{ijk} L_i L_j \bar{E}_k + \lambda'_{ijk} L_i Q_j \bar{D}_k + \frac{1}{2} \lambda''_{ijk} \bar{U}_i \bar{D}_j \bar{D}_k \, ,
    \label{eq:W_RPV}
\end{equation}
where $L_i$, $\bar E_i$, $Q_i$, $\bar U_i$, and $\bar D_i$ are chiral superfields with generation indices $i,j,k \in 
\{1,2,3\}$ and $H_u$ is one of the MSSM Higgs superfields. The $\lam_{ijk}$, $\lam'_{ijk}$ (and $\lam''_{ijk}$) are 
dimensionless Yukawa couplings, which imply lepton- (baryon-) number violating interactions and $\kappa_i$ are 
dimensionful bilinear couplings violating lepton number.
The Lagrangian in superfield-component form, as well as a complete list 
of RPV Feynman rules can be found, \textit{e.g.}, in Appendix L of Ref.~\cite{Dreiner:2008tw}. 

The RPV-MSSM allows for an unstable light long-lived neutralino, which we focus on here. A very light neutralino 
is necessarily dominantly bino-like \cite{Choudhury:1999tn} and a light bino currently avoids all experimental
and astrophysical constraints even if it is massless \cite{Gogoladze:2002xp, Dreiner:2003wh, Dreiner:2011fp}. In the 
following, we provide a concise overview of the production and decay of the lightest neutralino via $LL\bar E$ or 
$LQ\bar D$ operators. Moreover, for simplicity, we consider the lightest neutralino to be the LSP in this work. Note 
that we use the two-component fermion notation reviewed in Ref.~\cite{Dreiner:2008tw}.


\subsubsection{Neutralino production and decay via $LL\bar E$ operators}
For non-zero $LL \bar E$ couplings, neutralinos can be produced through charged lepton decays and can decay 
to lighter leptons. We provide the explicit general forms of the total decay widths for both the charged lepton and
neutralino decays, employing the matrix element and necessary phase space integration in 
Appendix ~\ref{sec:appendixA}. For all the relevant processes, we assume that all sfermions appear at the energy scale 
of the decaying particle as mass degenerate, $i.e.$, $m_{\tilde{f}} \simeq m_\text{SUSY} \gg \mchi, m_{\ell^{\pm}}$. 
Further, the R-parity conserving neutralino gauge coupling is $g'$ ($U(1)_Y$), since the neutralino is bino-like.
As a result, we can write out the coefficients appearing in the matrix elements as, (\textit{cf.}~Appendix~\ref{sec:appendixA})
\begin{equation}
    \label{eq:coeff_LLE}
  c_{ijk} \simeq -\frac{\frac{1}{\sqrt{2}} \lambda_{ijk} g' }{m_\text{SUSY}^2} \qquad \text{and} \qquad k_{ijk} \simeq \frac{\sqrt{2} \lambda_{ijk} g' }{m_\text{SUSY}^2}\, .
\end{equation}
Using Eq.~\eqref{eq:generic_LLE_width}, the relevant production widths can then be expressed as
\begin{align}
    \varGamma(\ell_k^\pm \rightarrow \tilde\chi_1^0 + \nu_i + \ell_j^\pm) = \varGamma_{LL\bar E} (\ell_k^\pm; \tilde\chi_1^0, \nu_i, \ell_j^\pm)[k_{ijk},c_{ijk},c_{ijk}] \, , \\
    \varGamma(\ell_k^\pm \rightarrow \tilde\chi_1^0 + \bar\nu_i + \ell_j^\pm) = \varGamma_{LL\bar E} (\ell_k^\pm; \tilde\chi_1^0, \bar\nu_i, \ell_j^\pm)[c_{ikj},k_{ikj},c_{ikj}] \, .
\end{align}
Similarly, the total widths of the subsequent neutralino decays $\tilde\chi_1^0 \rightarrow \overset{\scriptscriptstyle(-)}{\nu_i} + \ell_j^- + \ell_k^+$ can be written as
\begin{align}
    \varGamma(\tilde\chi_1^0 \rightarrow \nu_i + \ell_j^- + \ell_k^+) = \varGamma_{LL\bar E} (\tilde\chi_1^0; \nu_i, \ell_j^-, \ell_k^+)[c_{ijk},c_{ijk},k_{ijk}] \, , \\
    \varGamma(\tilde\chi_1^0 \rightarrow \bar\nu_i + \ell_j^- + \ell_k^+) = \varGamma_{LL\bar E} (\tilde\chi_1^0; \bar\nu_i, \ell_j^-, \ell_k^+)[c_{ikj},k_{ikj},c_{ikj}] \, .
\end{align}




\subsubsection{Neutralino production and decay via $LQ \bar D$ operators}
Via the $LQ \bar D$-operators mesons can decay into a bino accompanied by a lepton $l_i$. Subsequently, the bino 
decays via the same or another $LQ \bar D$-operator to a lepton and two quarks, where the latter again hadronize 
into a meson (for a light enough bino).
We consider charged mesons $M_{ab}^+$ with quark flavor content $(u_a \,\bar d_b)$ as well as neutral 
mesons $M_{ab}^0$ composed of $(d_a\, \bar d_b)$ and their charge conjugated equivalents. 
Neutral mesons composed of $(u_a \bar u_b)$ only contribute to higher multiplicity processes as $M \rightarrow \widetilde{\chi}^0+l_i+M^{\prime}$, where $M^{\prime}$ denotes a lighter meson and $l_i$ a charged lepton. We do \textit{not} consider them here, since these are phase space suppressed by two to three orders of magnitude~\cite{deVries:2015mfw}.
 Both the bino production and decay width are therefore given by Ref.~\cite{deVries:2015mfw}:
\begin{align}
    \varGamma(M_{ab} \rightarrow  \tilde\chi_1^0 + l_i) &= \frac{\lambda^{\frac{1}{2}}(m_{M_{ab}}^2, \mchi^2, m_{l_i}^2)}{64\pi m_{M_{ab}}^3} |G_{iab}^{S,f}|^2 (f_{M_{ab}}^S)^2\left(m_{M_{ab}}^2-\mchi^2-m_{l_i}^2\right) \, , \\
    \varGamma(\tilde\chi_1^0 \rightarrow M_{ab} + l_i) &= \frac{\lambda^{\frac{1}{2}}(\mchi^2, m_{M_{ab}}^2, m_{l_i}^2)}{128\pi \mchi^3} |G_{iab}^{S,f}|^2 (f_{M_{ab}}^S)^2\left(\mchi^2+m_{l_i}^2-m_{M_{ab}}^2\right) \, ,
\end{align}
where $\lambda^{\frac{1}{2}}(x,y,z)=\sqrt{x^2+y^2+z^2-2xy-2xz-2yz}\,$ is the square root of the K\"all\'en function and 
$l_i = \ell^{\pm}_i$ or $\nu_i$ depending on the charge of $M_{ab}$. The coefficients $|G_{iab}^{S,f}|^2$ include the 
trilinear RPV couplings and are defined in Ref.~\cite{deVries:2015mfw}, together with the meson scalar decay constants $f_{M_{ab}}^S$.

Furthermore, the $L_iL_j\bar E_j$ and $L_iQ_j\bar D_j$ operators additionally open the decay mode $\tilde
\chi_1^0 \rightarrow (\gamma + \nu_i, \gamma + \bar\nu_i)$ at the one-loop level, 
\textit{cf.}~Ref.~\cite{Dreiner:2022swd}, but this signature is not considered in the present paper since to our knowledge the corresponding HNL-decay has to-date not been searched for.

Besides the displaced-vertex signature related to the decay channels computed above, it is possible that the lightest 
neutralino is so long-lived that it does not decay inside the considered detector and appears as missing energy. 


\subsection{Heavy neutral leptons}
Heavy neutral leptons are a common feature of many SM extensions attempting to give an underlying explanation
of the observed neutrino sector. The simplest HNL model one can implement is 
\begin{equation}
    \lgr \supset i \hat N^\dagger_{\alpha} \bar\sigma^\mu \partial_\mu \hat N^{\alpha} - \left[\left(Y_
    \nu\right)^i_{\alpha} \left(\varPhi^0 \hat\nu_i \hat N^\alpha - \varPhi^+\ell_i \hat N^\alpha\right) + \frac{1}{2}  
    M^\alpha_\beta \hat N_\alpha \hat N^\beta  + \text{h.c.} \right] \, ,
\end{equation}
where  $i=1,2,3$, $\varPhi^{+}$ and $\varPhi^{0}$ are the components of the SM $SU(2)_L$ Higgs doublet, and $\ell_i$ 
are the charged lepton mass eigenstates. Fields with a hat, $\hat\nu$ and $\hat N$, are the states before mass-diagonalizing the neutral lepton sector. $\left(Y_\nu\right)^i_{\alpha}$ are dimensionless Yukawa couplings and $M_\beta^\alpha = 
\text{diag}(M_{\hat N_1}, ... )$ is a diagonal mass matrix.
The index $\alpha=1,2,3,\ldots$ labels the (arbitrary many) HNLs in the theory. During electroweak symmetry breaking, 
the Higgs obtains a vacuum expectation value (vev) $v/\sqrt{2}$ with $v=246$ GeV, which gives rise to mixing of the 
HNLs $\hat N_\alpha$ with active neutrinos $\hat\nu_i$, described by the mass matrix $M_{\nu N}$:
\begin{equation}
\label{eq:hnl_massmatrix}
    M_{\nu N} = \begin{pmatrix}
    \mathbb{0}_{3\times 3} && M_D \\
    M_D^T && M
    \end{pmatrix}\, .
\end{equation}
Here, the off-diagonal entries are given by $(M_D)^i_\alpha = (Y_\nu)^i_\alpha v / \sqrt{2}$. The mass matrix can be 
perturbatively Takagi-block-diagonalized by introducing a unitary matrix $U$ \cite{Dreiner:2008tw}. For simplicity, we 
assume there is only one kinematically relevant HNL in our study. In this case, the gauge eigenstate $\hat \nu_i$ 
receives first-order contributions from the mass eigenstate $N$, proportional to the following mixing matrix entry:
\begin{equation}
    \label{eq:hnl_u}
    U_i \equiv  U^i_4 \equiv (Y_\nu^*)^{i}_1 \frac{v}{\sqrt{2} M} \, .
\end{equation}
The interaction Lagrangian with the neutrino mass eigenstates $\nu_i$ and $N$ is then given by:
\begin{equation}
    \label{eq:hnl_lgr}
    \lgr \supset -\frac{g}{\sqrt{2}} U^i_4 W^-_\mu \ell_i^\dagger \bar\sigma^\mu N - \frac{g}{2 c_W} U^i_4 Z_\mu \nu_i^\dagger \bar\sigma^\mu N + \text{h.c.} \, ,
\end{equation}
where $g$ is the $SU(2)_L$ gauge coupling and $c_W = \cos\theta_W$ is the cosine of the weak mixing angle.






\subsection{The phenomenology connecting the light bino LSP and the HNL}
The phenomenologies of the RPV-MSSM with a light bino $\tilde{\chi}_1^0$ and the SM extensions with one relevant HNL 
 turn out to be very similar. This is not surprising, as the HNL and the bino have the same gauge quantum numbers after electroweak 
symmetry breaking. 
Currently existing bounds in the HNL parameter space spanned by $(m_N, U^i_4)$ can thus be 
translated into bounds in the light-bino-RPV parameter space $(m_{\tilde\chi_1^0}, \lam/m_\text{SUSY}^2)$, where 
$\lam$ labels here any appropriate $LL \bar E$ or $LQ \bar D$ coupling.

An additional analogy between the theories can be constructed by considering the bilinear RPV couplings $\kappa_i$~\cite{Diaz:1998vf}, see Eq.~\eqref{eq:W_RPV}.
After integrating out the heavy higgsinos in the neutral fermion sector of the RPV-MSSM, one obtains a tree-level mixing of neutrinos with the bino, which is of the form
\begin{equation}
    \label{eq:nu_bino_bilinear_mixing}
    \lgr \supset \frac{g'}{2} \left( v_i -\,\frac{v_d\kappa_i}{\kappa^0} \right)\hat \nu_i \tilde\chi_1^0 + \text{h.c.}\, ,
\end{equation}
where $v_i$ and $v_d$ are the vevs of the sneutrinos and the MSSM Higgs $H_d$, respectively. $\kappa^0$ is the Higgsino mass parameter and $g'$ the $U(1)_Y$ gauge coupling.
This mixing can be interpreted as the off-diagonal entries in the neutral lepton mass matrix given in Eq.~\eqref{eq:hnl_massmatrix}, such that the elements of the matrix $U$, \textit{cf.}~Eqs.~\eqref{eq:hnl_u} and \eqref{eq:hnl_lgr}, can be mapped to the neutralino-neutrino mixing considered here:
\begin{equation}
    U^i_4 = \frac{g'}{2 \mchi} \left( v_i -\,\frac{v_d\kappa_i}{\kappa^0} \right) \, .
\label{eq:bilinearmixingmatrix}
\end{equation}
For phenomenological computations, the Lagrangian given in Eq.~\eqref{eq:hnl_lgr} can be used, which would correspond to the trivial replacement
\begin{equation}
\begin{tikzpicture}[baseline=-0.1cm]
    \setlength{\feynhandlinesize}{1.2pt}
    \begin{feynhand}
        \vertex [particle] (f1) at (-1.5, 0) {$\hat \nu$};
        \vertex [particle] (f2) at (1.5, 0) {$N$};
        \propagator [plain, insertion={[size=4pt]0.5}] (f1) to (f2);
    \end{feynhand} 
\end{tikzpicture}
\qquad\longleftrightarrow \qquad
    \begin{tikzpicture}[baseline=-0.1cm]
    \setlength{\feynhandlinesize}{1.2pt}
    \begin{feynhand}
        \vertex [particle] (f1) at (-1.5, 0) {$\hat \nu$};
        \vertex [particle] (f2) at (1.5, 0) {$\tilde\chi_1^0$};
        \propagator [plain, insertion={[size=4pt]0.5}] (f1) to (f2);
    \end{feynhand} 
\end{tikzpicture}
\end{equation}
where the inserted crosses denote the mixing. Thus, current HNL exclusion limits can be directly translated into bounds 
on the mixing strength in Eq.~\eqref{eq:nu_bino_bilinear_mixing}.